\subsection*{C18. Complete measurement and positioning definitions}

\paragraph{Continuous floor-boundary distance.}
Let \(\mathcal{B}_P,\mathcal{B}_R\) contain midpoint samples at equal-length intervals on the spherical floor boundaries of layouts \(P,R\). Let \(\Gamma_P,\Gamma_R\) be their continuous shorter great-circle arcs between consecutive corners. With spherical angular distance \(\theta\) in degrees, the symmetric distance is
\begin{equation}
 \begin{aligned}
 2S_b(P,R)={}&\frac{1}{|\mathcal B_P|}\sum_{p\in\mathcal B_P}
 \min_{r\in\Gamma_R}\theta(p,r)\\
 &+\frac{1}{|\mathcal B_R|}\sum_{r\in\mathcal B_R}
 \min_{p\in\Gamma_P}\theta(r,p).
 \end{aligned}
 \label{eq:supp-boundary-distance}
\end{equation}
Every sample is compared with the other layout's continuous boundary. The implementation uses 2,048 samples in each direction and gives the two directed arithmetic means equal weight.

\paragraph{Global and local height.}
Heights interpolate linearly along floor edges. For cyclic corner heights \(h_i\) and floor-edge lengths \(\ell_i\), the perimeter-weighted mean height and relative mean-height error are
\begin{equation}
 \mu_P=\frac{\sum_i\ell_i(h_i+h_{i+1})/2}{\sum_i\ell_i},
 \qquad H_m(P,R)=\frac{|\mu_P-\mu_R|}{\mu_R}.
 \label{eq:supp-height-mean}
\end{equation}
Subtract each layout's own mean to obtain \(\widetilde h_P(x)=h_P(x)-\mu_P\) on \(\partial A_P\), and similarly for \(R\). For midpoint samples \(\mathcal X_P\) at equal floor-perimeter intervals, define the nearest horizontal-floor positions
\begin{equation}
 \mathcal N_R(x)=\underset{y\in\partial A_R}{\arg\min}\,\|x-y\|_2.
 \label{eq:supp-nearest-floor-set}
\end{equation}
The directed centered-height error is
\begin{equation}
 \operatorname{MSE}^{c}_{P\rightarrow R}
 =\frac{1}{|\mathcal X_P|}\sum_{x\in\mathcal X_P}
 \max_{y\in\mathcal N_R(x)}
 [\widetilde h_P(x)-\widetilde h_R(y)]^2.
 \label{eq:supp-height-directed}
\end{equation}
The maximum resolves ties in nearest floor position. Exchanging \(P,R\) gives the reverse error, and
\begin{equation}
 H_c(P,R)=\frac1{\mu_R}
 \sqrt{\frac{\operatorname{MSE}^{c}_{P\rightarrow R}
 +\operatorname{MSE}^{c}_{R\rightarrow P}}2}.
 \label{eq:supp-height-center}
\end{equation}
The calculation uses 4,096 samples per direction. Both height errors are dimensionless; a changed corner can affect the global and local components.

\paragraph{Score scales and geometric comparisons.}
The boundary-loss denominator is \(32\,\mathrm{deg}^2\), the direction scale is \(8^\circ\), and the direction loss approaches 1.5. Global and local height scales are 0.8 and 0.25. Thus \(S_b=\sqrt{32}^\circ\) contributes unit boundary loss, an \(8^\circ\) direction residual contributes 0.75, and each height loss equals one at its stated scale. The L2 combination lies between the maximum and sum of the floor and direction losses. Q reaches 100 when the footprint and height profiles agree with the reference and the assessed layout has zero internal direction residual.

The controlled examples distinguish local height, structural completeness and geometric deviation. In the square-room control with reference height 2.6 camera heights, raising one side by 20\% and lowering the opposite side by 20\% preserves the footprint and mean height. Prism IoU remains 1, \(H_m=0\), \(H_c=0.1633\), and Q falls from 100 to 70.09. Column-wise layout-depth RMSE is 0.4437\,m under the 1.6\,m camera-height convention. In the 3\,m controls, raising one corner by 2.4\,m or every corner by 0.6\,m gives the same corner mean and prism IoU of 83.33\%; Q is 46.97 and 94.12, respectively. Omitting the deepest tested narrow recess removes four turns, with prism IoU at 97.09\% and Q at 96.67. Turn counts record this structural loss alongside geometric scores (Supplement A).

A real team-endpoint comparison also separates overlap, structure and boundary accuracy. Floor IoU changes from 0.8414 to 0.8451 and Q from 75.03 to 78.80. Bottom-corner recall at \(9^\circ\) rises from 0.5 to 1, while layout-depth RMSE falls from 0.2359\,m to 0.1851\,m. These component measurements describe simultaneous changes in the complete geometry.

\paragraph{Support rule.}
For candidate group \(G\) in team \(S\), distinct-person support is
\begin{equation}
 s(G)=|\{w\in S:\exists\,a_{wj}\in G\}|.
 \label{eq:supp-support}
\end{equation}
For fixed \(G\), the binary decision \(z\in\{0,1\}\) incurs disagreement
\begin{equation}
 \mathcal L_G(z)=s(G)(1-z)+[|S|-s(G)]z.
 \label{eq:supp-support-disagreement}
\end{equation}
The MV50 rule \(z=\mathrm{1}[2s(G)\geq|S|]\) minimizes this disagreement, retaining the corner when the two alternatives tie.

\paragraph{Periodic position.}
The unwrapping anchor is an observed pair minimizing the sum of the larger endpoint distances:
\begin{equation}
 a^*=\underset{a\in G}{\arg\min}\sum_{b\in G}
 \max\{\theta(p_a^t,p_b^t),\theta(p_a^b,p_b^b)\}.
 \label{eq:supp-position-anchor}
\end{equation}
Let \(x_0\) be its aligned horizontal coordinate, \(W=1024\) the panorama width, and \(x_j\) the aligned source coordinates. With \(\operatorname{wrap}_W(d)=((d+W/2)\bmod W)-W/2\), the retained corner coordinates are
\begin{align}
 x_G&=\left[x_0+\operatorname{median}_j\{\operatorname{wrap}_W(x_j-x_0)\}\right]\bmod W,\label{eq:supp-median-x}\\
 y_G^t&=\operatorname{median}_j y_j^t,\qquad y_G^b=\operatorname{median}_j y_j^b.\label{eq:supp-median-y}
\end{align}
All three medians use the same supporting people with equal weights. The spread must fit within an open half-panorama; C16 specifies this check, tie resolution and the source-ring connection procedure.
